\chapter{Lukisan Geometri}\label{ch:g5:geometry}

Penggaris, penggaris siku, jangka: dengan tiga alat ini dan sedikit
cara kerja, setiap gambar yang dijelaskan dengan kata-kata dapat
dibangun dengan tepat. Bab ini melukis \omterm{def:g4:geometry:def}{garis tegak lurus}, \omterm{def:g4:geometry:def}{garis
sejajar}, segitiga dan \omterm{def:g4:geometry:polygon}{segi empat} istimewa --- keterampilan dulu,
teori belakangan (\cref{ch:g6:lines,ch:g6:shapes}).

\section{Garis tegak lurus dan sejajar melalui sebuah titik}

\begin{method}[Dua lukisan dasar]\label{met:g5:geometry:basic}
\emph{\omterm{def:g4:geometry:def}{Garis tegak lurus} terhadap $d$ melalui $A$}: geserlah
penggaris siku sepanjang $d$ sampai tepi \omterm{def:g3:shapes:rightangle}{sudut siku-sikunya} yang
lain mencapai $A$; gambar; lalu tandai \omterm{def:g3:shapes:rightangle}{sudut siku-sikunya}.

\emph{\omterm{def:g4:geometry:def}{Garis sejajar} dengan $d$ melalui $A$}: tekan penggaris biasa
pada penggaris siku sebagai rel yang diletakkan sepanjang $d$; geser
penggaris siku itu sampai ke $A$; lalu gambar. (Periksa: jarak
antara kedua garis itu sama di mana-mana.)
\end{method}

\begin{example}[Merangkai lukisan]\label{ex:g5:geometry:chain}
Banyak gambar adalah rangkaian dari kedua gerakan itu. Persegi
panjang $ABCD$ dengan $AB = 5$ cm dan $BC = 3$ cm: gambarlah $[AB]$;
\omterm{def:g4:geometry:def}{garis tegak lurus} di $A$; \omterm{def:g4:geometry:def}{garis tegak lurus} di $B$; tandai $D$ dan
$C$ pada jarak $3$ cm di sisi yang sama; lalu hubungkan. Setiap
langkahnya adalah satu lukisan dasar.
\end{example}

\section{Segitiga dengan jangka}

\begin{method}[Segitiga dari tiga sisi]\label{met:g5:geometry:triangle}
Untuk membangun segitiga dengan sisi $6$ cm, $5$ cm, $4$ cm:
\begin{enumerate}
  \item gambarlah sisi yang terpanjang, $[AB]$ dengan $AB = 6$ cm;
  \item busur berpusat $A$ dengan \omterm{def:g3:shapes:shapes}{jari-jari} $5$ cm;
  \item busur berpusat $B$ dengan \omterm{def:g3:shapes:shapes}{jari-jari} $4$ cm;
  \item perpotongannya adalah \omterm{def:g4:geometry:polygon}{titik sudut} ketiga $C$; lalu hubungkan.
\end{enumerate}
Jangka menjamin panjangnya: setiap titik pada busur yang pertama
berjarak tepat $5$ cm dari $A$.
\end{method}

\begin{omfigure}
\begin{tikzpicture}[scale=0.95]
  \coordinate (A) at (0,0);
  \coordinate (B) at (6*0.72,0);
  \coordinate (C) at (2.7,2.376);
  \draw[thick] (A) -- (B);
  \draw[thick, omDef] (A) -- (C);
  \draw[thick, omThm] (B) -- (C);
  \draw[omDef, thin] (C) arc[start angle=41.3, delta angle=20,
    radius=3.6];
  \draw[omDef, thin] (C) arc[start angle=41.3, delta angle=-20,
    radius=3.6];
  \draw[omThm, thin] (C) arc[start angle=124.4, delta angle=20,
    radius=2.88];
  \draw[omThm, thin] (C) arc[start angle=124.4, delta angle=-20,
    radius=2.88];
  \fill (A) circle (1.5pt) node[below left, font=\small] {$A$};
  \fill (B) circle (1.5pt) node[below right, font=\small] {$B$};
  \fill (C) circle (1.5pt) node[above, font=\small] {$C$};
  \node[below, font=\small] at (2.16,0) {$6$ cm};
  \node[omDef, font=\small] at (0.85,1.45) {$5$ cm};
  \node[omThm, font=\small] at (3.9,1.35) {$4$ cm};
\end{tikzpicture}

{\small Dua busur, satu perpotongan: segitiga $6$--$5$--$4$ yang
dibangun dengan jangka. (Kalau kedua sisi pendeknya bersama-sama
lebih pendek daripada alasnya, busurnya tidak akan bertemu --- lebih
lanjut pada \cref{ch:g7:triangles}.)}
\end{omfigure}

\begin{example}[Sama kaki dan sama sisi secara cuma-cuma]\label{ex:g5:geometry:special}
Bukaan jangka yang sama untuk kedua busur $\Rightarrow$ segitiga
sama kaki. Bukaan yang sama dengan alasnya $\Rightarrow$ segitiga
sama sisi. Jangka tidak mengukur: ia \emph{menyalin} panjang, dan
itu malah lebih baik.
\end{example}

\section{Segi empat istimewa, lewat lukisan}

\begin{method}[Membangun keluarganya]\label{met:g5:geometry:quadrilaterals}
\begin{enumerate}
  \item \emph{Persegi panjang}: dua pasang \omterm{def:g4:geometry:def}{garis sejajar} yang
        berpotongan \omterm{def:g4:geometry:def}{tegak lurus} --- atau rangkaian pada
        \cref{ex:g5:geometry:chain};
  \item \emph{persegi}: persegi panjang yang keempat sisinya memakai
        panjang yang sama --- bawalah panjang itu dengan jangka;
  \item \emph{belah ketupat}: empat sisi sama panjang --- gambarlah
        satu sisi, lalu dua busur jangka dari kedua ujungnya dengan
        bukaan yang sama menetapkan dua \omterm{def:g4:geometry:polygon}{titik sudut} \omterm{def:g3:division:remainder}{sisanya} (belah
        ketupat adalah dua segitiga sama kaki yang direkatkan
        sepanjang sebuah diagonal).
\end{enumerate}
Setelah membangun, selalu \emph{periksalah} dengan alatnya:
penggaris siku pada pojoknya, jangka pada sisinya.
\end{method}

\begin{example}[Mengenali dari sebuah gambaran]\label{ex:g5:geometry:recognize}
``Sebuah \omterm{def:g4:geometry:polygon}{segi empat} dengan empat sisi $4$ cm dan satu \omterm{def:g3:shapes:rightangle}{sudut
siku-siku}'' --- bangunlah: keempat sisi yang sama panjang memaksa
bentuk belah ketupat, dan menekan satu pojoknya menjadi \omterm{def:g3:shapes:rightangle}{sudut
siku-siku} meluruskan semua pojok yang lain: hasilnya adalah persegi.
Ada gambaran yang diam-diam menjelaskan bangun yang lebih istimewa
daripada yang diumumkannya. (Mengapa \omterm{def:g3:shapes:rightangle}{sudut siku-sikunya} menular ke
keempat pojok dijelaskan pada \cref{ch:g7:central}.)
\end{example}

\section{Lingkaran dalam lukisan}

\begin{example}[Lingkaran sebagai alat]\label{ex:g5:geometry:circle}
``Carilah semua titik yang berjarak $3$ cm dari $A$'': itulah
lingkaran berpusat $A$ dengan \omterm{def:g3:shapes:shapes}{jari-jari} $3$ cm. ``Carilah semua
titik yang berjarak $3$ cm dari $A$ \emph{dan} $4$ cm dari $B$'':
gambarlah kedua lingkarannya; perpotongannya (dua, satu, atau \omterm{def:g1:counting:zero}{nol}
titik) adalah jawabannya. Peta harta karun dan lukisan segitiga
memakai gagasan yang persis sama.
\end{example}

\section{Latihan}

\begin{exercise}[$\star$]\label{exo:g5:geometry:1}
Gambarlah sebuah garis $d$ dan titik $A$ berjarak kira-kira $3$ cm
darinya. Lukislah \omterm{def:g4:geometry:def}{garis tegak lurus} terhadap $d$ melalui $A$, lalu
\omterm{def:g4:geometry:def}{garis sejajar} dengan $d$ melalui $A$. Tandailah \omterm{def:g3:shapes:rightangle}{sudut siku-sikunya}.
\end{exercise}

\begin{exercise}[$\star$]\label{exo:g5:geometry:2}
Lukislah persegi panjang $ABCD$ dengan $AB = 7$ cm dan $BC = 4$ cm,
mengikuti \cref{ex:g5:geometry:chain}. Periksalah sudut keempatnya
dengan penggaris siku: kalau lukisannya tepat, sudut itu otomatis
siku-siku.
\end{exercise}

\begin{exercise}[$\star$]\label{exo:g5:geometry:3}
Lukislah segitiga dengan sisi $7$ cm, $6$ cm dan $3$ cm
(\cref{met:g5:geometry:triangle}).
\end{exercise}

\begin{exercise}[$\star$]\label{exo:g5:geometry:4}
Lukislah segitiga sama kaki dengan alas $5$ cm dan kaki $7$ cm; lalu
segitiga sama sisi dengan sisi $6$ cm.
\end{exercise}

\begin{exercise}[$\star$]\label{exo:g5:geometry:5}
Lukislah persegi dengan sisi $5$ cm. Alat mana yang menjamin \omterm{def:g3:shapes:rightangle}{sudut
siku-sikunya}? Alat mana yang menjamin sisinya sama panjang?
\end{exercise}

\begin{exercise}[$\star$]\label{exo:g5:geometry:6}
Lukislah belah ketupat dengan sisi $4$ cm yang bukan persegi
(pilihlah sendiri bukaan untuk pojok pertamanya). Periksalah keempat
sisinya dengan jangka.
\end{exercise}

\begin{exercise}[$\star$]\label{exo:g5:geometry:7}
Gambarlah dua titik $A$ dan $B$ dengan $AB = 5$ cm. Lukislah semua
titik yang berjarak $4$ cm dari $A$ dan $3$ cm dari $B$. Ada berapa
titik seperti itu?
\end{exercise}

\begin{exercise}[$\star$]\label{exo:g5:geometry:8}
Tulislah sebuah program melukis (langkah bernomor, titik yang diberi
nama) untuk persegi panjang $6$ cm $\times$ $2$ cm yang di atasnya
ada segitiga sama kaki dengan kaki $4$ cm (bentuk pensil). Mintalah
teman sekelasmu mengerjakannya.
\end{exercise}

\begin{exercise}[$\star\star$]\label{exo:g5:geometry:9}
Cobalah melukis segitiga dengan sisi $8$ cm, $3$ cm dan $4$ cm. Apa
yang terjadi pada busurnya? Jelaskan dalam satu kalimat mengapa
segitiga itu tidak mungkin ada.
\end{exercise}

\begin{exercise}[$\star\star$]\label{exo:g5:geometry:10}
Lukislah belah ketupat yang diagonalnya $6$ cm dan $4$ cm, dengan
mengetahui bahwa diagonal belah ketupat berpotongan \omterm{def:g4:geometry:def}{tegak lurus} di
tengah-tengahnya: gambarlah diagonalnya \emph{lebih dulu}, lalu
hubungkan keempat ujungnya.
\end{exercise}

\begin{exercise}[$\star\star$]\label{exo:g5:geometry:11}
Seekor kambing diikat pada pancang $A$ dengan tali $4$ m, di
lapangan datar yang punya pagar lurus melintas $3$ m dari $A$.
Gambarlah denahnya (1 cm untuk 1 meter): daerah mana yang dapat
dicapai kambing itu? Di mana pagarnya membatasi dia?

\end{exercise}
